%% kapa: exam=1 round=20 year=2009 period=2 form=A subject=appraisal_law subject_ko=부동산관계법규 seq=32 number=72
%% source: Q-Net 감정평가사 기출문제 (cst00309) artlSeq=5024178 file=extracted/1008/5024178/2009년 제20회 감정평가사 1차시험 기출문제/2교시.hwp
%% conversion: kordoc 4.24.51 -> tools/gen1.py (auto); review=proofread-sonnet; flags=none
\begin{question}{72}
건축법령상 이행강제금에 관한 설명으로 옳은 것은?
\begin{choices}
  \item 허가권자는 계고 없이 이행강제금을 부과ㆍ징수할 수 있다.
  \item 위법이 시정된 경우에는 이미 부과된 이행강제금은 징수하지 아니한다.
  \item 이행강제금의 부과 및 징수절차는 국고금관리법 시행규칙을 준용한다.
  \item 이행강제금은 1년에 1회 이내의 범위에서 그 시정명령이 이행될 때까지 반복하여 부과ㆍ징수할 수 있다.
  \item 긴급한 경우에는 시정명령 없이 이행강제금을 부과할 수 있다.
\end{choices}
\end{question}
\answer{3}
